% GENERATED, edit the French unit or the translation table
\mxunit{premiere-c08}{Random variables and dispersion}\label{C08}
\begin{keypoint}[Chapter path]
\textbf{Objectives.} Build a distribution; interpret events; calculate expectation, variance, standard deviation; use linearity; simulate samples.

\textbf{Prerequisites.} Probability, weighted mean, squares, trees and lists.

\textbf{Yearly reference point.} 19--28 January: 8 hours\end{keypoint}
\section*{Diagnostic: six questions, 10 minutes}\label{C08-D}
Answer without consulting the lesson. Mark the questions for which help is needed. This diagnostic is not graded.
\begin{enumerate}
\item Calculate the mean of 2, 4, 6.
\item Check whether $(\dfrac{1}{2},\dfrac{1}{3},\dfrac{1}{6})$ can be a distribution.
\item Calculate $(2-5)^2$.
\item If $P(A)=0.6$, calculate $P(\overline A)$.
\item A gross gain of 8 units and a stake of 3 give what net gain?
\item Calculate $\sqrt{0.25}$.
\end{enumerate}
\paragraph{After the diagnostic.} Revisit the flagged topics with the teacher; redo a similar question, then explain the rule used.
\paragraph{Speaking and logic.} Explain why a theoretical mean can be an impossible gain. Present the difference between a distribution, a simulation and a data set.
\mxnewpage
\section*{Investigate before formalising}\label{C08-RECH}
Create two games with three equally likely outcomes, both with expectation 0 but different standard deviations. Check their distributions and explain the difference.
\paragraph{Organisation.} Individual investigation 5 min, comparison in pairs 5 min, class discussion 5 min. Keep a record of one attempt, even an unfinished one. A conjecture must then be justified.
\section{Variable, event and distribution}\label{C08-U1}
\begin{keypoint}[Definition and results]
On a finite sample space $\Omega$, a real random variable X is a function that associates a real number with each outcome. The event $\{X=a\}$ groups the outcomes giving the value a. Its distribution associates with each distinct value $x_i$ the probability $p_i=P(X=x_i)$, with $p_i\geq0$ and $\sum p_i=1$.
\end{keypoint}
\paragraph{Status and justification.} Definitions. The events associated with distinct values form a partition of the sample space; this justifies the sum of the probabilities.
\begin{example}A fair die gives a gain X equal to 2 for 1 or 2, and to $-1$ otherwise. The distribution is $P(X=2)=\dfrac{1}{3}$, $P(X=-1)=\dfrac{2}{3}$. $\{X\leq0\}=\{3;4;5;6\}$: probability $\dfrac{2}{3}$.\end{example}\label{C08-EX1}
\mxinclude{figures/C08-U1}
\begin{pitfall}The values are not always equally likely even if the outcomes are. Group the outcomes that give the same gain.\end{pitfall}
\begin{application}Two fair coins are tossed. X counts the heads. Write its distribution and $P(X\leq1)$.\end{application}\label{C08-AP1}
\section{Expectation and linearity}\label{C08-U2}
\begin{keypoint}[Definition and results]
The expectation is $E(X)=\sum_i x_ip_i$. It describes the long-term theoretical mean, not a result that is necessarily possible on a single trial. For variables X and Y on the same finite sample space and real numbers a,b: $E(aX+b)=aE(X)+b$ and $E(X+Y)=E(X)+E(Y)$, with no assumption of independence.
\end{keypoint}
\paragraph{Status and justification.} Summing over the outcomes $\omega$, $E(X+Y)=\sum_\omega(X(\omega)+Y(\omega))P(\omega)$; distributing the sum gives the result. The constant b has expectation b because the probabilities sum to 1.
\begin{example}For the previous distribution $P(X=2)=\dfrac{1}{3}$, $P(X=-1)=\dfrac{2}{3}$, $E(X)=0$: the game is fair if X is the net gain. For $Y=3X-2$, $E(Y)=3\times0-2=-2$.\end{example}\label{C08-EX2}
\begin{pitfall}An expectation is not an individual guarantee; a game can be losing at every observed horizon despite a positive expectation.\end{pitfall}
\begin{application}A gross gain is 10 with probability $\dfrac{1}{5}$ and 0 otherwise. What stake makes the game fair?\end{application}\label{C08-AP2}
\section{Variance, standard deviation and König--Huygens}\label{C08-U3}
\begin{keypoint}[Definition and results]
Writing $\mu=E(X)$, $V(X)=E((X-\mu)^2)=\sum_i(x_i-\mu)^2p_i$. The standard deviation is $\sigma(X)=\sqrt{V(X)}$. The König--Huygens formula gives $V(X)=E(X^2)-E(X)^2$. A variance is always positive or zero; the standard deviation has the same unit as X, the variance a squared unit.
\end{keypoint}
\paragraph{Status and justification.} Expand $(X-\mu)^2=X^2-2\mu X+\mu^2$, apply linearity and $E(X)=\mu$: $E(X^2)-2\mu^2+\mu^2$. We obtain the formula.
\begin{example}For X with values 0 and 4 with probabilities $\dfrac{3}{4}$ and $\dfrac{1}{4}$, $E(X)=1$, $E(X^2)=4$, $V(X)=3$ and $\sigma=\sqrt3$. Direct calculation: $(0-1)^2(\dfrac{3}{4})+(4-1)^2(\dfrac{1}{4})=3$.\end{example}\label{C08-EX3}
\begin{pitfall}$E(X^2)$ is not $E(X)^2$. A variance calculated as negative signals an error or a very small rounding error to be dealt with explicitly.\end{pitfall}
\begin{application}For X with values $-1,1$ equally likely, calculate the expectation, variance and standard deviation.\end{application}\label{C08-AP3}
\section{Transforming and comparing gains}\label{C08-U4}
\begin{keypoint}[Definition and results]
For real a,b, $V(aX+b)=a^2V(X)$ and $\sigma(aX+b)=|a|\sigma(X)$. A translation moves the mean without changing the dispersion; a multiplication changes the scale. Two games with the same expectation can have very different dispersions.
\end{keypoint}
\paragraph{Status and justification.} $aX+b-E(aX+b)=a(X-E(X))$. Squaring and then taking the expectation gives $a^2V(X)$. The square root of $a^2$ is $|a|$.
\begin{example}A game A always pays 2. A game B pays 0 or 4 with equal probability. Both have expectation 2; variances 0 and 4. Their distributions are therefore not interchangeable despite the same mean.\end{example}\label{C08-EX4}
\begin{pitfall}The variance of a sum is not, in general, the sum of the variances; no rule of this type is used without hypotheses.\end{pitfall}
\begin{application}$E(X)=3$, $V(X)=4$. Give the characteristics of $Y=-2X+5$.\end{application}\label{C08-AP4}
\section{Samples and observed mean}\label{C08-U5}
\begin{keypoint}[Definition and results]
A sample of size n is here a list of n independent realisations of the same variable X. Its observed mean m varies from one sample to another. We compare m with $\mu=E(X)$; repetitions allow us to observe the proportion of samples satisfying $|m-\mu|\leq \dfrac{2\sigma}{\sqrt n}$. This threshold is an experimental task of the programme, not a claim of an exact coverage rate.
\end{keypoint}
\paragraph{Status and justification.} Experimentation: neither a convergence theorem nor the normal distribution is proved here. We check the size of the lists, the calculation of the mean and the distinction between n and the number N of samples.
\begin{example}For X with values 0 and 2 equally likely, $\mu=1$, $\sigma=1$. For n=100, threshold $\dfrac{2}{10}=0.2$: we count the means in $[0.8;1.2]$. A result of 956 out of 1000 means an observed frequency of 0.956 for this simulation only.\end{example}\label{C08-EX5}
\begin{pitfall}Do not copy a single draw n times; do not announce that the proportion is necessarily 95 \%.\end{pitfall}
\begin{application}For $\mu=5$, $\sigma=3$ and n=36, give the experimental interval asked for.\end{application}\label{C08-AP5}
\section*{Exercise path}\label{C08-PARC}
Start with exercises 1 and 2. Then choose one of exercises 3 or 4 and do it without help. Exercises 5 to 7 are a reserve for transfer, investigation or programming: the teacher designates one. Exercise 8 is revisited after a delay.
\begingroup
\mxexercisehead{C08-E1}{A distribution to build}
A fair six-sided die, with faces numbered 1 to 6, is rolled. The random variable $X$ associates with the result the remainder of its Euclidean division by 3. For example, a result divisible by 3 gives $X=0$.
\begin{enumerate}[label=\arabic*.]
\item For each of the six faces, state the value taken by $X$.
\item Present the probability distribution of $X$ in a table and check that the probabilities sum to 1.
\item Calculate $P(X<2)$, stating which values of $X$ are retained.
\end{enumerate}
\mxexercisehead{C08-E2}{Two variance calculations}
The random variable $X$ takes the values $-2$, 0 and 4 with respective probabilities $\dfrac14$, $\dfrac12$ and $\dfrac14$.
\begin{enumerate}[label=\arabic*.]
\item Calculate its expectation $\mu=E(X)$ and the second moment $E(X^2)$.
\item Calculate $V(X)$ using the König--Huygens formula, then deduce the exact standard deviation.
\item Recover the variance from the definition $V(X)=E((X-\mu)^2)$. Present the three contributions to the calculation.
\end{enumerate}
\mxexercisehead{C08-E3}{A fair game}
A wheel has five sectors of equal angle, bearing the numbers 0, 0, 0, 5 and 10. The five sectors are assumed equally likely. The number obtained represents a gross gain $G$, in monetary units. To play, one pays a stake $m$; the net gain is therefore $X=G-m$.
\begin{enumerate}[label=\arabic*.]
\item Establish the distribution of $G$, then determine the stake $m$ that makes the game fair, that is, such that $E(X)=0$.
\item After studying variance, calculate the standard deviation of the gross gain. Deduce that of the net gain when the stake is fair.
\item Explain why changing the stake changes the expectation of the net gain but not its standard deviation.
\end{enumerate}
\mxinclude{figures/C08-E3}
\mxexercisehead{C08-E4}{Same mean, different risk}
Two fictitious investments are offered. The first yields a certain gain $X=1$. The second yields a gain $Y=-2$ with probability $\dfrac12$ and a gain $Y=4$ with probability $\dfrac12$. The gains are expressed in the same monetary unit.
\begin{enumerate}[label=\arabic*.]
\item Compare the expectations of the two gains, then their variances.
\item Someone states: ``These investments are equivalent since they have the same mean.'' Explain what this sentence omits and what the variance makes it possible to distinguish here.
\end{enumerate}
\mxexercisehead{C08-E5}{Linearity without independence}
A fair die numbered 1 to 6 is rolled. Let $X$ be the result of the roll and define $Y=7-X$. The two variables are therefore obtained from the same roll.
\begin{enumerate}[label=\arabic*.]
\item Calculate $E(X)$ and $E(Y)$, then $E(X+Y)$. Check the linearity of expectation on this example.
\item Are the variables $X$ and $Y$ independent? Justify using, for example, the events $\{X=1\}$ and $\{Y=6\}$.
\item Does the linearity of expectation require independence? Conclude using your results.
\end{enumerate}
\mxexercisehead{C08-E6}{A minimal loss}
A random variable $X$ takes finitely many values and has an expectation $\mu$. For a real number $a$, the mean squared error associated with the constant forecast $a$ is measured by $E((X-a)^2)$.
\begin{enumerate}[label=\arabic*.]
\item Prove that $E((X-a)^2)=V(X)+(a-\mu)^2$. You may write $X-a=(X-\mu)+(\mu-a)$.
\item Deduce the value of $a$ that minimises this error, then the minimum value. Justify the uniqueness of the optimal choice.
\end{enumerate}
\mxexercisehead{C08-E7}{Letter frequencies}
We wish to measure the frequency of the letter ``a'' in a text. The denominator is the number of alphabetic characters: spaces, digits and punctuation marks are excluded. Capital letters are treated as lower case; accented letters remain distinct from ``a''.
\begin{enumerate}[label=\arabic*.]
\item Write a Python function that returns this frequency. State the result chosen when the text contains no letters.
\item Test the function with \code{Aba !}, then with \code{123 !} and with the empty string. Explain the expected results before running the program.
\end{enumerate}
\mxexercisehead{C08-E8}{Delayed review}
An experiment whose probability of success is $\dfrac12$ is repeated twice, independently. The random variable $X$ counts the number of successes obtained.
\begin{enumerate}[label=\arabic*.]
\item Recover the distribution of $X$ using the four possible sequences of results.
\item Calculate $E(X)$ and $V(X)$ starting from this distribution.
\end{enumerate}
\paragraph{Homework.} Do only the task designated at the end of the session, in the dated progression: 10--20 min at most. The exercises depend on the topics already studied, so some questions are deferred. Revisit exercise 8 one week after the chapter (10 min). If the allotted time is exceeded, note the difficulty and bring your attempts; do not pile up the whole reserve.
\endgroup
\mxneed{12}\section*{To do in the CQFD textbook}
These exercises are done in your textbook. The progression indicates the session; they replace the homework planned for that day (10 to 20 minutes in total).
\paragraph{CQFD, exercise 28 p.~306, all questions.}\label{C08-M1} Calculate the expectation from the distribution table; keep the probabilities associated with each value.
\paragraph{CQFD, exercise 31 p.~306, all questions; use the distribution of exercise 28.}\label{C08-M2} Calculate the variance then the standard deviation, respecting the rounding to the nearest tenth asked for by the textbook.
\begingroup\small
\mxnewpage\section*{MCQ 1: understand and choose}\label{C08-Q1}
Five single-answer questions for practice. Choose one answer per question; then explain one distractor. This chapter format does not reproduce the official EAM format.
\begin{enumerate}
\item A probability distribution has a sum of probabilities equal to
\begin{enumerate}[label=\Alph*.,itemsep=2pt]
\item $0$
\item $1$
\item $n$
\item the mean
\end{enumerate}
\item X is 0 or 4 with probabilities $\dfrac{3}{4}$ and $\dfrac{1}{4}$; $E(X)=$
\begin{enumerate}[label=\Alph*.,itemsep=2pt]
\item $1$
\item $2$
\item $3$
\item $4$
\end{enumerate}
\item $V(X)$ equals
\begin{enumerate}[label=\Alph*.,itemsep=2pt]
\item $E(X^2)$
\item $E(X)^2$
\item $E(X^2)+E(X)^2$
\item $E(X^2)-E(X)^2$
\end{enumerate}
\item If $V(X)=9$, $\sigma(X)=$
\begin{enumerate}[label=\Alph*.,itemsep=2pt]
\item $81$
\item $9$
\item $3$
\item $\dfrac{1}{3}$
\end{enumerate}
\item $E(2X+1)$ equals
\begin{enumerate}[label=\Alph*.,itemsep=2pt]
\item $2E(X)$
\item $2E(X)+1$
\item $E(X)+1$
\item $E(X)^2+1$
\end{enumerate}
\end{enumerate}
\endgroup
\begingroup\small
\mxneed{18}\section*{MCQ 2: understand and choose}\label{C08-Q2}
Five single-answer questions for practice. Choose one answer per question; then explain one distractor. This chapter format does not reproduce the official EAM format.
\begin{enumerate}
\item $V(X+5)$ equals
\begin{enumerate}[label=\Alph*.,itemsep=2pt]
\item $V(X)+5$
\item $V(X)+25$
\item $V(X)$
\item $5V(X)$
\end{enumerate}
\item $\sigma(-3X)$ equals
\begin{enumerate}[label=\Alph*.,itemsep=2pt]
\item $3\sigma(X)$
\item $-3\sigma(X)$
\item $9\sigma(X)$
\item $\sigma(X)$
\end{enumerate}
\item The expectation is necessarily
\begin{enumerate}[label=\Alph*.,itemsep=2pt]
\item a possible value
\item an integer
\item positive
\item none of these three statements
\end{enumerate}
\item $E(X+Y)=E(X)+E(Y)$ requires
\begin{enumerate}[label=\Alph*.,itemsep=2pt]
\item independence
\item no independence
\item X=Y
\item X positive
\end{enumerate}
\item For $\sigma=2$, n=100, $\dfrac{2\sigma}{\sqrt n}=$
\begin{enumerate}[label=\Alph*.,itemsep=2pt]
\item $4$
\item $0.04$
\item $0.4$
\item $2$
\end{enumerate}
\end{enumerate}
\subsection*{True or false? Justify.}\label{C08-VF}
\begin{enumerate}
\item A variance can be negative.
\item Two variables with the same expectation have the same distribution.
\item Adding a constant to a variable does not change its standard deviation.
\end{enumerate}
\begin{keypoint}[Self-assessment]I build a complete distribution; I distinguish gross from net; I check both variance formulas; I do not confuse frequency with probability.

For each objective, indicate: achieved without help; achieved with help; to be reworked. Write one example and one mistake to avoid.\end{keypoint}
\endgroup
\ifmxteacher\mxinclude{corriges/C08}\fi
